\chapter{Catalytic holonomy, excursion geometry, and resource reuse}\label{ch:holonomy}
\index{holonomy}
\index{Schreier graph}
\index{catalytic conjugation}
\index{dirty workspace}
\index{resource amortization}


\section{Subgroup metrics and Schreier loops}
Assume first that the restricted and ambient models are groups
\[
H=\langle\mathcal S_H\rangle
\le
G=\langle\mathcal S_G\rangle.
\]
For $h\in H$, the intrinsic and ambient word metrics are
\[
\ell_H(h)=|h|_{\mathcal S_H},
\qquad
\ell_G(h)=|h|_{\mathcal S_G}.
\]
The quotient geometry is the right Schreier graph $H\backslash G$.

\begin{theorem}[Holonomy-compression identity]
Let every generator have unit cost.  For $h\in H$, $\ell_G(h)$ is the minimum length of a based loop in $H\backslash G$ whose lifted word has holonomy $h$.  The intrinsic length $\ell_H(h)$ is the minimum length among loops that remain in the base fiber.
\end{theorem}
\begin{proof}
An ambient word $s_k\cdots s_1=h$ projects to the loop
\[
H,
Hs_1,
Hs_2s_1,
\ldots,
Hh=H.
\]
Conversely, a based Schreier loop has product in $H$.  Restricting the path to the base fiber forces every prefix to lie in $H$, hence every edge label can be replaced by a restricted generator under the chosen intrinsic Cayley structure.
\end{proof}

Thus
\[
\Delta_{H\hookrightarrow G}(h)
=
\text{intrinsic fiber cost}
-
\text{shortest holonomy loop cost}.
\]
This is the geometric content of a dirty-workspace shortcut.

\begin{figure}[ht]
\centering
\includegraphics[width=.90\textwidth]{catalytic_excursion_general.pdf}
\caption{A catalytic excursion leaves a restricted fiber, reuses an ambient state, and returns with nontrivial holonomy.}
\label{fig:catalytic-excursion}
\end{figure}

\section{The basic conjugation inequality}
Let $g\in G$ and $k\in H$ satisfy $gkg^{-1}\in H$.  Then
\[
\ell_G(gkg^{-1})
\le 2\ell_G(g)+\ell_H(k).
\]
Define the intrinsic conjugation stretch
\[
\operatorname{Stretch}_g(k)
=
\ell_H(gkg^{-1})-\ell_H(k).
\]

\begin{proposition}[Catalytic stretch criterion]
If
\[
\operatorname{Stretch}_g(k)>2\ell_G(g),
\]
then $gkg^{-1}$ has positive defect.
\end{proposition}
The opening and closing cost is paid once.  Every useful consumer inside the excursion contributes to the stretch.

\section{Many consumers and catalytic excess}
Suppose the middle word is
\[
k=h_m\cdots h_1,
\qquad h_i\in H.
\]
The ambient word
\[
gh_m\cdots h_1g^{-1}
\]
has cost at most
\[
2c(g)+\sum_i c(h_i).
\]
If the restricted images $gh_ig^{-1}$ require costs $r_i$ and a lower-bound invariant proves that their obligations cannot be shared, then
\[
\ell_H(gkg^{-1})\ge\sum_i r_i.
\]
A sufficient inequality is
\[
\sum_i(r_i-c(h_i))>2c(g).
\]
The left side is the gross catalytic benefit.  The right side is the restoration cost.

\begin{definition}[Catalytic excess]
For a normal form with $r$ matched excursions and $C$ useful cross interactions, define
\[
\mathcal E=C-2r
\]
when each consumer contributes one unit and each opening or closing costs one unit.
\end{definition}
Positive excess is necessary in the simplest transvection model, but it is not sufficient.  An all-cross configuration may admit a mirror word of equal restricted length.  Temporal asymmetry or phase holonomy is also required.

\section{Excursion budgets}
Counting ambient-only gates is only one measure of how far a circuit leaves its restricted model.
\begin{definition}[Excursion profile]
For a circuit path $X_0,X_1,\ldots,X_L$, an excursion profile may include
\begin{enumerate}
\item the number $q$ of ambient-only gates,
\item the maximum Schreier distance $d(HX_t,H)$,
\item the rank of a forbidden block,
\item the bandwidth of a forbidden dependency,
\item the number of simultaneously open catalytic intervals,
\item the intersection graph of those intervals,
\item the maximum algebraic degree created temporarily,
\item the maximum ancilla dimension or branch count.
\end{enumerate}
\end{definition}
The established linear family shows that a rank-one, bandwidth-one excursion using exactly two forbidden gates can save linearly many operations.  Thus geometric excursion size need not control computational benefit.

\section{Catalytic interval geometry}
Associate to every matched opening and closing operation an interval
\[
I_j=[t_j^{\mathrm{open}},t_j^{\mathrm{close}}].
\]
The family of intervals can be disjoint, nested, laminar, or crossing.

\begin{principle}[First higher-order frontier]
Independent or disjoint intervals often reduce to additive copies of a one-catalyst analysis.  Crossing intervals are the first place where one dirty resource can alter another and where a multiplicative distortion larger than the clean one-carrier value may occur.
\end{principle}

Define the catalytic width
\[
w_{\mathrm{cat}}=
\max_t
\#\{j:t\in I_j\}.
\]
One may then study
\[
D_n^{(w)}
=
\max_{U}
\left(
\ellRes(U)-\ellAmb^{(w)}(U)
\right),
\]
where the ambient circuit has width at most $w$.  The unitriangular snapshot family gives $D_n^{(1)}=\Omega(n)$ in its corresponding interval model.

\section{Four recurrent catalytic mechanisms}
\subsection{Basis-change catalysis}
An ambient gate $g$ changes the frame in which several restricted operations are cheap:
\[
h_i\mapsto gh_ig^{-1}.
\]
Hadamard and Fourier transforms are the basic quantum examples.  The frame is restored at the end.

\subsection{Commutator catalysis}
An ambient commutator creates a new direction:
\[
[g,h]=ghg^{-1}h^{-1}.
\]
For continuous controls,
\[
e^{\varepsilon X}e^{\varepsilon Y}e^{-\varepsilon X}e^{-\varepsilon Y}
=
e^{\varepsilon^2[X,Y]+O(\varepsilon^3)}.
\]
The desired effective generator may lie in the restricted algebra even though the path does not.

\subsection{Degree catalysis}
A Toffoli, CCZ, or non-Gaussian gate can create higher algebraic degree or higher moments temporarily.  If the nonlinear object is reused and canceled, the final target may remain affine, diagonal, or Gaussian.

\subsection{Branch and workspace catalysis}
Ancillas, measurements, reset, teleportation, and classical feedforward enlarge the realization category.  The resource is a temporary branch, subsystem, or borrowed state rather than a group element.

\section{When defects are impossible}
A useful positive theorem is a criterion guaranteeing geodesic convexity.
\begin{proposition}[One-Lipschitz retraction criterion]
Suppose there is a retraction $\pi:G\to H$ with $\pi|_H=\operatorname{id}_H$ such that every ambient generator $s\in\mathcal S_G$ satisfies
\[
\ell_H(\pi(s))\le c(s).
\]
Then
\[
\ell_H(h)\le\ell_G(h)
\]
for every $h\in H$, hence the two metrics agree on $H$.
\end{proposition}
\begin{proof}
Apply $\pi$ gate by gate to an ambient word for $h$.  The projected factors form a restricted word of no greater cost.
\end{proof}

The unitriangular counterexample proves that no such generator-wise one-Lipschitz retraction exists for the lower transvection model inside the all-transvection model.

\section{Weighted restoration and phase transitions}
Let every restricted gate cost one and every ambient-only gate cost $\lambda$.  For a catalyst used by $v$ consumers with unit benefit, the displayed construction has net gain at least
\[
v-2\lambda.
\]
The critical weight is therefore target and mechanism dependent.

\begin{definition}[Critical catalyst penalty]
For a target family $U_n$, define
\[
\lambda_c(U_n)
=
\sup\{\lambda:\ell_{\lambda,\mathrm{amb}}(U_n)<\ellRes(U_n)\}.
\]
\end{definition}
Determining $\lambda_c$ exactly is a refined synthesis problem.  It can be more physically meaningful than unweighted gate count.

\section{Restoration is the source of nontriviality}
If the ambient resource need not be restored, the opening cost can dominate the comparison and many examples become trivial.  Catalytic escape is strongest when the following all hold:
\begin{enumerate}
\item the final target lies entirely in the restricted class,
\item every auxiliary subsystem or violated constraint is restored,
\item the ambient-only resource is used more than once,
\item the restricted lower bound charges the same repeated obligations separately,
\item the gain survives a realistic weighted cost.
\end{enumerate}

\begin{conjecturebox}
In many finite gate hierarchies, minimal defects should admit a bounded-width catalytic normal form.  Proving or refuting this would turn an unrestricted synthesis search into a search over a finite family of excursion motifs.
\end{conjecturebox}
